% Báo cáo độc lập, tập trung vào kết quả; không cần file dữ liệu hoặc .bib.
% Biên dịch hai lần bằng pdfLaTeX, XeLaTeX hoặc LuaLaTeX.
% Số liệu đã nhận tới 07/10/2026; biên tập rút gọn ngày 08/10/2026.
\documentclass[11pt,a4paper]{article}
\usepackage{iftex}
\ifPDFTeX
  \usepackage[T5]{fontenc}
  \usepackage[utf8]{inputenc}
\else
  \usepackage{fontspec}
  \setmainfont{TeX Gyre Termes}
  \setsansfont{TeX Gyre Heros}
  \setmonofont{Latin Modern Mono}
\fi
\usepackage[vietnamese]{babel}
\usepackage[margin=2.2cm]{geometry}
\usepackage{amsmath,amssymb,booktabs,array,longtable,tabularx}
\usepackage{xcolor,xurl}
\usepackage[unicode,colorlinks=true,linkcolor=blue!55!black,
  citecolor=blue!55!black,urlcolor=blue!55!black]{hyperref}
\setlength{\parindent}{0pt}
\setlength{\parskip}{5pt}
\setlength{\emergencystretch}{3em}
\renewcommand{\arraystretch}{1.15}
\newcommand{\code}[1]{\texttt{\detokenize{#1}}}
\newcommand{\ACC}{\operatorname{ACC}}
\newcommand{\MAP}{\operatorname{MAP}}
\newcommand{\NR}{\textnormal{NR}}
\newcommand{\SKIP}{\textnormal{SKIP}}
\title{Báo cáo thực nghiệm sơ bộ\\
  Tái triển khai FlashSinkhorn và đánh giá FlashOPW}
\author{Dự án \texttt{flash-opw}}
\date{Kết quả ghi nhận đến ngày 07 tháng 10 năm 2026}

\begin{document}
\maketitle
\begin{abstract}
Báo cáo trình bày kết quả tám benchmark synthetic FlashSinkhorn trên RTX 5080
và các nhóm đánh giá FlashOPW. FlashSinkhorn là bản tự triển khai theo paper;
FlashOPW dùng công thức affine/Taylor của tài liệu main, còn đối chứng OPW
journal giữ công thức và score riêng. Nhóm 1 có bằng chứng correctness và
hội tụ trên pilot; nhóm 2 hoàn tất ablation; nhóm 3 hoàn tất tuning riêng
từng phương pháp trên TRAIN. Chưa có kết quả official TEST đầy đủ và
benchmark hiệu năng GPU cuối cùng cho nhóm 4--6 trong dữ liệu đã nhận.
\end{abstract}
\tableofcontents
\newpage

\section{Phạm vi và môi trường}
Các solver FlashSinkhorn, FlashOPW và baseline OPW journal là bản tự triển
khai, không phải code chính thức của tác giả. Kết quả hỗ trợ tính đúng trên
các trường hợp đã kiểm tra; chưa xác nhận tái hiện đầy đủ thực nghiệm gốc.
Benchmark GPU dùng RTX 5080 khoảng 16 GB; paper FlashSinkhorn dùng A100 80 GB.
Các run OPW đã kiểm tra dùng PyTorch 2.11.0+cu128, Triton 3.6.0;
reference CPU dùng FP64, Flash và dense GPU dùng FP32.

Số liệu FlashSinkhorn được đọc từ log mean, chưa có raw repeats và môi trường
đầy đủ của full run. Kết quả OPW nhóm 2--3 đã được kiểm tra lại từ các file
kết quả; nhóm 1 GPU có log parity và bảng hội tụ, còn thiếu bộ file tổng kết
để xác nhận numerical parity toàn matrix. Trạng thái nhóm 4--6 phản ánh
thông tin đã nhận tới ngày 07/10/2026.

\section{FlashSinkhorn: tám benchmark và đối chiếu paper}
\label{sec:fs}
\subsection{Cấu hình chung}
\begin{tabularx}{\linewidth}{@{}lX@{}}
\toprule
Thành phần & Cấu hình chính \\
\midrule
Input & $X,Y\sim U([0,1]^d)$, $n=m$, uniform weights,
  squared Euclidean cost, $\varepsilon=0.1$. \\
Forward / F+B & 10 vòng cố định; TF32; symmetric và alternating;
  10 warmup, 50/30 repeats. F+B gồm forward và backward. \\
HVP & FP32; solve 100 vòng ngoài timing, CG 50 bước,
  damping $10^{-5}$; 10 warmup, 20 repeats. \\
Đo thời gian & Torch/KeOps dùng CUDA events; JAX dùng synchronized wall clock.
  Tensorized precompute cost ngoài forward timing. \\
Đo bộ nhớ & Peak PyTorch allocated sau warmup, gồm live tensors;
  không trừ bộ nhớ nền và không đo tổng VRAM. \\
\bottomrule
\end{tabularx}

Cấu hình toán học bám sát Appendix H.2 \cite{flashpaper}. JAX HVP dùng công
thức OTT-Hessian với adapter CG, nên chưa đại diện timing implementation
Lineax nguyên bản. Paper cũng mô tả số repeats HVP không nhất quán giữa
H.2 và captions Tables 15--16. Validation GPU đã ghi nhận 133 tests đạt;
benchmark full có 208 số đo, 12 memory skip và 8 ô HVP không được báo cáo.

\textbf{Cách đọc bảng:} ô speedup $a/b$ ghi lần lượt RTX/paper, không phải
phép chia hai số. Speedup là $t_{\rm baseline}/t_{\rm Flash}$; lớn hơn 1
nghĩa là Flash nhanh hơn. KeOps/Tensorized forward và F+B chia cho Flash-sym,
JAX chia cho Flash-alt; HVP chia cho Flash-sym. Ô in đậm đánh dấu đảo thứ tự
nhanh/chậm so paper. \SKIP\ là bỏ trước theo memory estimate; \NR\ là không
được báo cáo; NC là chưa có số paper tương ứng trong nguồn bảng được dùng.
OOM/OOT là trạng thái của paper, không tự gán cho server. Số đo baseline
đầy đủ nằm ở Phụ lục~\ref{app:fs-tables}.

\subsection{Experiment 1: forward theo số điểm, $d=64$}
$n=m=5000,10000,20000,30000,40000,50000$; đối chiếu Tables 8/10/12.
% BEGIN FS COMPARE forward_n
\begin{center}\footnotesize
\setlength{\tabcolsep}{5pt}
\begin{tabular}{@{}rrrrrr@{}}
\toprule
$n=m$ & \shortstack{Flash-sym\\ms} & \shortstack{Flash-alt\\ms} & \shortstack{KeOps speedup\\RTX/Paper} & \shortstack{Tensorized speedup\\RTX/Paper} & \shortstack{JAX speedup\\RTX/Paper} \\
\midrule
5000 & 2.259 & 2.182 & 18.98/14.70 & 11.11/3.30 & 4.03/1.90 \\
10000 & 6.613 & 7.64 & 11.64/15.40 & 15.88/6.60 & 3.30/2.50 \\
20000 & 24.337 & 26.502 & 11.22/9.40 & 17.34/8.10 & 3.12/3.60 \\
30000 & 49.996 & 54.45 & 8.08/11.00 & \SKIP/OOM & 3.26/2.80 \\
40000 & 88.538 & 89.333 & 8.99/8.30 & \SKIP/OOM & 3.44/3.60 \\
50000 & 135.652 & 138.441 & 9.67/9.90 & \SKIP/OOM & 3.53/4.20 \\
\bottomrule
\end{tabular}
\end{center}
% END FS COMPARE forward_n
Flash nhanh hơn mọi baseline có số đo trong sweep, cùng chiều với paper.
Tensorized tại $n\geq30000$ bị preflight skip; paper ghi OOM.

\subsection{Experiment 2: forward theo dimension, $n=m=20000$}
$d=4,8,16,32,64,128,256,512,1024$; đối chiếu Tables 8/10/12.
% BEGIN FS COMPARE forward_d
\begin{center}\footnotesize
\setlength{\tabcolsep}{5pt}
\begin{tabular}{@{}rrrrrr@{}}
\toprule
$d$ & \shortstack{Flash-sym\\ms} & \shortstack{Flash-alt\\ms} & \shortstack{KeOps speedup\\RTX/Paper} & \shortstack{Tensorized speedup\\RTX/Paper} & \shortstack{JAX speedup\\RTX/Paper} \\
\midrule
4 & 13.477 & 11.842 & 2.27/1.60 & 31.32/9.90 & 5.95/3.20 \\
8 & 13.455 & 11.806 & 3.30/2.60 & 31.37/9.90 & 6.07/3.30 \\
16 & 13.545 & 11.81 & 5.67/4.20 & 31.16/12.30 & 6.20/3.60 \\
32 & 13.508 & 11.85 & 10.63/9.70 & 31.25/10.50 & 6.44/4.30 \\
64 & 24.503 & 26.387 & 11.08/9.40 & 17.23/8.10 & 3.12/3.60 \\
128 & 45.02 & 47.333 & 9.72/10.80 & 9.38/4.60 & 2.19/3.00 \\
256 & 120.886 & 126.721 & 10.06/19.30 & 3.49/3.00 & 1.13/2.40 \\
512 & 238.024 & 240.36 & 9.89/17.60 & 1.77/1.70 & \textbf{0.96/1.60} \\
1024 & 468.222 & 473.183 & 9.97/14.40 & 0.90/0.70 & 1.01/1.20 \\
\bottomrule
\end{tabular}
\end{center}
% END FS COMPARE forward_d
JAX nhanh hơn Flash-alt tại $d=512$, trái chiều paper. Tại $d=1024$,
Tensorized nhanh hơn Flash-sym ở cả hai phép đo.

\subsection{Experiment 3: forward cộng backward theo số điểm, $d=64$}
Cùng sáu kích thước của experiment 1; đối chiếu Tables 9/11/13.
% BEGIN FS COMPARE backward_n
\begin{center}\footnotesize
\setlength{\tabcolsep}{5pt}
\begin{tabular}{@{}rrrrrr@{}}
\toprule
$n=m$ & \shortstack{Flash-sym\\ms} & \shortstack{Flash-alt\\ms} & \shortstack{KeOps speedup\\RTX/Paper} & \shortstack{Tensorized speedup\\RTX/Paper} & \shortstack{JAX speedup\\RTX/Paper} \\
\midrule
5000 & 3.835 & 3.254 & 13.25/15.00 & 8.32/3.60 & 2.82/2.70 \\
10000 & 8.21 & 9.212 & 11.00/12.30 & 16.11/5.60 & 3.16/2.90 \\
20000 & 28.502 & 30.628 & 11.07/8.50 & 18.62/7.90 & 3.07/3.60 \\
30000 & 58.154 & 62.554 & 8.08/10.30 & \SKIP/OOM & 3.26/3.70 \\
40000 & 101.63 & 102.717 & 9.02/8.90 & \SKIP/OOM & 3.48/4.20 \\
50000 & 155.901 & 160.041 & 9.68/10.50 & \SKIP/OOM & 4.03/4.00 \\
\bottomrule
\end{tabular}
\end{center}
% END FS COMPARE backward_n
Flash giữ lợi thế trước các baseline có số đo, cùng chiều với paper.

\subsection{Experiment 4: forward cộng backward theo dimension, $n=m=20000$}
Cùng chín dimension của experiment 2; đối chiếu Tables 9/11/13.
% BEGIN FS COMPARE backward_d
\begin{center}\footnotesize
\setlength{\tabcolsep}{5pt}
\begin{tabular}{@{}rrrrrr@{}}
\toprule
$d$ & \shortstack{Flash-sym\\ms} & \shortstack{Flash-alt\\ms} & \shortstack{KeOps speedup\\RTX/Paper} & \shortstack{Tensorized speedup\\RTX/Paper} & \shortstack{JAX speedup\\RTX/Paper} \\
\midrule
4 & 15.632 & 13.957 & 2.31/1.50 & 33.92/10.60 & 5.82/3.70 \\
8 & 15.596 & 13.926 & 3.31/2.60 & 33.99/10.70 & 5.75/3.90 \\
16 & 15.573 & 13.92 & 5.67/4.10 & 34.04/12.80 & 5.83/4.80 \\
32 & 15.602 & 14.097 & 10.58/9.70 & 33.99/11.50 & 6.05/4.70 \\
64 & 28.051 & 30.061 & 11.10/8.50 & 18.91/7.90 & 3.07/3.60 \\
128 & 51.968 & 54.245 & 11.15/11.50 & 10.22/4.90 & 2.11/2.70 \\
256 & 135.281 & 141.095 & 19.84/33.30 & 3.95/3.10 & 1.16/2.20 \\
512 & 273.385 & 275.735 & 62.52/OOT & 1.98/1.40 & \textbf{0.95/1.80} \\
1024 & 538.959 & 544.009 & 132.49/OOT & \textbf{1.03/0.90} & \textbf{0.97/1.10} \\
\bottomrule
\end{tabular}
\end{center}
% END FS COMPARE backward_d
JAX nhanh hơn Flash-alt tại $d=512,1024$, trái chiều paper.
Tensorized/Flash-sym tại $d=1024$ cũng đảo chiều quanh speedup 1;
chênh nhỏ cần raw samples để xác nhận độ ổn định.

\subsection{Experiment 5: bộ nhớ forward theo số điểm, $d=1024$}
Cùng sáu kích thước của experiment 1; đối chiếu Figure 3.
% BEGIN FS COMPARE memory_forward
\begin{center}\footnotesize
\setlength{\tabcolsep}{5pt}
\begin{tabular}{@{}rrrr@{}}
\toprule
$n=m$ & Flash-alt MB & Tensorized MB & Tensorized/Flash \\
\midrule
5000 & 95.314 & 428.307 & 4.49 \\
10000 & 174.08 & 1370.148 & 7.87 \\
20000 & 325.714 & 5187.666 & 15.93 \\
30000 & 482.607 & \SKIP & -- \\
40000 & 634.388 & \SKIP & -- \\
50000 & 787.202 & \SKIP & -- \\
\bottomrule
\end{tabular}
\end{center}
% END FS COMPARE memory_forward
Log-log slope là 0.918 cho Flash và 1.799 cho Tensorized. Xu hướng tăng
gần tuyến tính của Flash phù hợp paper; chưa có bảng MB paper cùng setting
để đối chiếu chính xác từng điểm.

\subsection{Experiment 6: bộ nhớ forward cộng backward, $d=1024$}
Cùng sweep và cách đo experiment 5; đối chiếu Figure 3.
% BEGIN FS COMPARE memory_backward
\begin{center}\footnotesize
\setlength{\tabcolsep}{5pt}
\begin{tabular}{@{}rrrr@{}}
\toprule
$n=m$ & Flash-alt MB & Tensorized MB & Tensorized/Flash \\
\midrule
5000 & 101.008 & 562.464 & 5.57 \\
10000 & 184.974 & 2104.03 & 11.37 \\
20000 & 345.042 & 8182.158 & 23.71 \\
30000 & 512.45 & \SKIP & -- \\
40000 & 673.041 & \SKIP & -- \\
50000 & 839.925 & \SKIP & -- \\
\bottomrule
\end{tabular}
\end{center}
% END FS COMPARE memory_backward
Slope là 0.920 cho Flash và 1.931 cho Tensorized. Tại $n=20000$,
peak allocation của Tensorized bằng khoảng 23.71 lần Flash.

\subsection{Experiment 7: HVP theo số điểm, $d=64$}
$n=m=5000,6000,7000,8000,9000,10000,20000,30000,40000,50000$;
đối chiếu Tables 15/16 và Section 4.1.
% BEGIN FS COMPARE hvp_n
\begin{center}\footnotesize
\setlength{\tabcolsep}{5pt}
\begin{tabular}{@{}rrrrrr@{}}
\toprule
$n=m$ & \shortstack{Flash-sym\\ms} & \shortstack{KeOps\\ms} & \shortstack{JAX\\ms} & \shortstack{KeOps speedup\\RTX/Paper} & \shortstack{JAX speedup\\RTX/Paper} \\
\midrule
5000 & 65.31 & 193.927 & 290.295 & 2.97/6.10 & 4.44/3.70 \\
6000 & 100.726 & 216.518 & 409.974 & 2.15/6.20 & 4.07/4.50 \\
7000 & 127.179 & 237.857 & 566.703 & 1.87/4.20 & 4.46/3.60 \\
8000 & 153.503 & 264.787 & 698.838 & 1.72/4.20 & 4.55/3.90 \\
9000 & 199.077 & 292.444 & 1056.555 & 1.47/4.10 & 5.31/4.20 \\
10000 & 236.438 & 315.049 & 1289.244 & 1.33/4.00 & 5.45/4.60 \\
20000 & 912.952 & 931.432 & \NR & 1.02/NC & \NR/NC \\
30000 & 2032.442 & 1454.346 & \NR & 0.72/NC & \NR/NC \\
40000 & 3532.364 & 2745.119 & \NR & 0.78/NC & \NR/NC \\
50000 & 5337.679 & 4412.484 & \NR & \textbf{0.83/3.45}$^{*}$ & \NR/NC \\
\bottomrule
\end{tabular}
\end{center}
% END FS COMPARE hvp_n
KeOps nhanh hơn Flash tại $n\geq30000$ trong reproduction. Tại $n=50000$,
quan hệ đảo chiều so paper; ô có dấu $^{*}$ dùng tỷ số 14.5 s/4.2 s
của Section 4.1. JAX $n>10000$ không được báo cáo trong run.

\subsection{Experiment 8: HVP theo dimension, $n=m=10000$}
$d=4,8,16,32,64,128,256,512$; đối chiếu Tables 15/16.
% BEGIN FS COMPARE hvp_d
\begin{center}\footnotesize
\setlength{\tabcolsep}{5pt}
\begin{tabular}{@{}rrrrrr@{}}
\toprule
$d$ & \shortstack{Flash-sym\\ms} & \shortstack{KeOps\\ms} & \shortstack{JAX\\ms} & \shortstack{KeOps speedup\\RTX/Paper} & \shortstack{JAX speedup\\RTX/Paper} \\
\midrule
4 & 77.987 & 124.908 & 151.746 & \textbf{1.60/0.70} & 1.95/1.50 \\
8 & 78.009 & 133.39 & 169.775 & 1.71/1.00 & 2.18/1.60 \\
16 & 78.061 & 159.214 & 230.08 & 2.04/2.00 & 2.95/2.20 \\
32 & 78.56 & 207.22 & 450.829 & 2.64/3.80 & 5.74/3.60 \\
64 & 236.456 & 314.645 & 1289.259 & 1.33/4.00 & 5.45/4.60 \\
128 & 619.745 & 761.873 & 4528.709 & 1.23/17.00 & 7.31/7.00 \\
256 & 972.212 & \NR & \NR & \NR/OOT & \NR/OOT \\
512 & 2107.318 & \NR & \NR & \NR/OOT & \NR/OOT \\
\bottomrule
\end{tabular}
\end{center}
% END FS COMPARE hvp_d
Tại $d=4$, Flash nhanh hơn KeOps, trái chiều paper. KeOps/Flash tại $d=128$
chỉ khoảng 1.23, thấp hơn nhiều so paper. Baseline $d=256,512$ không được
chạy trong phạm vi runner, không phải OOT thực đo trên RTX.

\subsection{Runtime tuyệt đối và kết luận}
Đối chiếu cùng shape với Table 18 cho forward và Section 4.1 cho HVP.
RTX và A100 khác hardware; không dựng runtime baseline paper bằng cách ghép
speedup và runtime từ những phép đo khác nhau.
% BEGIN FS COMPARE absolute
\begin{center}\footnotesize
\setlength{\tabcolsep}{5pt}
\begin{tabular}{@{}rrlrrl@{}}
\toprule
$n=m$ & $d$ & Method & RTX ms & Paper ms & Paper source \\
\midrule
10000 & 64 & flash-sym & 6.613 & 14.0 & Table 18 \\
10000 & 64 & flash-alt & 7.64 & 17.2 & Table 18 \\
20000 & 64 & flash-sym & 24.337 & 37.4 & Table 18 \\
20000 & 64 & flash-alt & 26.502 & 39.6 & Table 18 \\
50000 & 64 & flash-sym & 135.652 & 188.1 & Table 18 \\
50000 & 64 & flash-alt & 138.441 & 192.2 & Table 18 \\
20000 & 1024 & flash-sym & 468.222 & 377.5 & Table 18 \\
20000 & 1024 & flash-alt & 473.183 & 303.7 & Table 18 \\
50000 & 64 & flash-sym & 5337.679 & 4200 & Section 4.1 \\
50000 & 64 & keops & 4412.484 & 14500 & Section 4.1 \\
\bottomrule
\end{tabular}
\end{center}
% END FS COMPARE absolute
Reproduction giữ nhiều lợi thế forward/F+B và memory scaling, nhưng chưa
tái hiện toàn bộ thứ hạng và speedup paper. Chưa đủ bằng chứng xác định
riêng ảnh hưởng của GPU, baseline implementation và cơ chế đo.

\section{FlashOPW: phương pháp và cấu hình chung}
\label{sec:opw-method}
Đặt $D_{ij}=\|x_i-y_j\|^2$, $R_{ij}=(i/N-j/M)^2$,
$\ell^2(i,j)=R_{ij}/(N^{-2}+M^{-2})$. Hai công thức \cite{mainpdf,journalpdf}:
\begin{align}
C^{\rm main}_{ij}&=D_{ij}+\mu R_{ij}+q_0,
&\mu&=\lambda_1+\frac{\lambda_2}{2\sigma^2},\\
C^{\rm journal}_{ij}&=D_{ij}-\frac{\lambda_1}{1+R_{ij}}
+\lambda_2\left[\frac{\ell^2(i,j)}{2\sigma^2}
+\log(\sigma\sqrt{2\pi})\right],
&q_0&=\lambda_2\log(\sigma\sqrt{2\pi})-\lambda_1.
\end{align}
$\mu$ là hệ số suy ra, không phải siêu tham số độc lập. Main dùng score
Eq.19 $\langle a,f\rangle+\langle b,g\rangle-\lambda_2+q_0$;
journal dùng $\langle P,D\rangle$. Dense affine là reference của main,
không phải OPW journal. Hai định nghĩa prior có đơn vị sigma khác nhau.

Dữ liệu chính: FacesUCR, TRAIN 200, TEST 2050, 14 lớp,
chuỗi dài 131, $d=1$; giữ giá trị và độ dài gốc, uniform weights,
cost scale 1. Nhóm 1--3 dùng cùng input làm tròn FP32 cho reference CPU
FP64 và GPU FP32. Nhóm accuracy cập nhật $f$ rồi $g$.
Gallery là tập mẫu đối chiếu, query là mẫu cần dự đoán.
ACC@k dùng majority vote từ k hàng xóm; MAP tính toàn ranking,
coi các mẫu cùng lớp là relevant. NN agreement đo sự đồng ý giữa hai backend,
không phải accuracy theo nhãn thật.

\section{Các pilot FlashOPW ban đầu}
\subsection{Parity CPU/GPU và pilot k-NN}
Cùng 32 TRAIN gallery, 16 TEST query, 512 cặp; tham số
$(\lambda_1,\lambda_2,\sigma)=(1,0.1,1)$.
CPU pilot so main/affine ở 20 vòng; GPU pilot so ở 200 vòng với FP64 reference.
Journal OPW dùng 20 vòng, các entropic baseline còn lại dùng 100 vòng.
\begin{center}\small
\begin{tabular}{@{}llrr@{}}
\toprule
Pilot & Phương pháp & MAP\% & ACC@1\% \\
\midrule
CPU & Main / dense affine, 20 vòng & 68.329 & 62.50 \\
CPU & TLp & 82.482 & 81.25 \\
CPU & OPW / OPW-KL journal & 84.034 & 93.75 \\
GPU & Main / dense affine reference, 200 vòng & 68.449 & 62.50 \\
\bottomrule
\end{tabular}
\end{center}
Main/affine GPU có max score error $9.58348\cdot10^{-8}$ so FP64;
ranking và predictions trùng. Distance-matrix time Flash 13.1206 s,
dense CPU 101.7103 s: tỷ số 7.75 khác backend/precision và chưa kiểm soát tải,
không coi là same-GPU speedup. Các fixed-iteration pilot chưa thống nhất
mức hội tụ, nên không dùng để kết luận ưu thế cuối cùng giữa các phương pháp.

\subsection{Đối chứng Taylor, prior và score trên pilot}
CPU FP64 trên cùng 512 cặp, $(1,0.1,1)$; thay từng yếu tố:
\begin{center}\small
\begin{tabular}{@{}llllrr@{}}
\toprule
Order & Prior & Score & Vòng & MAP\% & ACC@1\% \\
\midrule
Taylor & relative & Eq.19 & 200 & 68.449 & 62.50 \\
Exact & relative & cùng quy ước potentials & 200 & 69.576 & 62.50 \\
Exact & journal & cùng quy ước potentials & 200 & 78.199 & 87.50 \\
Exact & journal & $\langle P,D\rangle$ & 200 & 83.375 & 87.50 \\
Exact & journal & $\langle P,D\rangle$ & 20 & 84.034 & 93.75 \\
\bottomrule
\end{tabular}
\end{center}
Chênh lệch phụ thuộc đường đối chứng, không phải phân rã nhân quả duy nhất.
Pilot TEST đã được xem trong phát triển; không gọi toàn quá trình là blind TEST.

\subsection{Tuning Flash-only ban đầu}
TRAIN holdout 28 gallery/14 query, thử ba hệ số suy ra $\mu=1.05,50,430$,
giữ $\lambda_1=1,\lambda_2=0.1$, 200 vòng; chọn ACC@1 rồi MAP.
Bộ được chọn có $\sigma\approx0.0319438$ ($\mu=50$).
Trên subset TEST 32 gallery/28 query, kết quả cuối của lượt fixed 200 vòng:
\begin{center}\small
\begin{tabular}{@{}lrrr@{}}
\toprule
Cấu hình & MAP\% & ACC@1\% & Max residual \\
\midrule
Default & 59.854 & 53.571 & $6.432\cdot10^{-4}$ \\
Selected & 81.986 & 92.857 & $1.185\cdot10^{-2}$ \\
\bottomrule
\end{tabular}
\end{center}
Selected cải thiện accuracy trên subset nhưng chưa đạt residual $10^{-3}$
cho mọi cặp. Kết quả này không thay tuning riêng từng phương pháp ở nhóm 3.

\section{Nhóm 1: correctness và hội tụ}
\label{sec:g1}
\textbf{Setting chính:} TRAIN pilot 16 gallery/14 query, 224 cặp;
default $(1,0.1,1)$ và selected Flash $(1,0.1,0.0319438)$.
So Flash IEEE FP32 với explicit-cost CPU FP64 ở cùng số vòng;
kiểm tra score, spatial cost, $P$ và $P@[1,Y]$ trên sáu shape synthetic
và các cặp TRAIN. Ngưỡng relative L2 của $P$/apply là $5\cdot10^{-3}$.

\textbf{Kết quả cuối:} CPU 52/52 parity checks đạt; log GPU 40/40 đạt.
Max relative L2 của $P$ GPU là $3.26\cdot10^{-6}$ với default và
$4.14\cdot10^{-5}$ với selected. Cả 12 score matrices có NN agreement
100\% với FP64 tại cùng số vòng. Kết quả hội tụ GPU cuối lượt 2000 vòng:
\begin{center}\small
\begin{tabular}{@{}lrrr@{}}
\toprule
Cấu hình & Max residual & Cặp đạt $10^{-3}$ & Cặp đạt $10^{-4}$ \\
\midrule
Default & $1.92\cdot10^{-5}$ & 224/224 & 224/224 \\
Selected Flash & $7.32\cdot10^{-5}$ & 224/224 & 224/224 \\
\bottomrule
\end{tabular}
\end{center}
Residual là $\max(\|P\mathbf1-a\|_1,\|P^\top\mathbf1-b\|_1)$.
Kết quả hỗ trợ correctness và đạt hai ngưỡng trên pilot;
NN agreement không chứng minh ACC 100\% hay numerical parity toàn matrix.
Còn thiếu file tổng kết GPU để đóng kiểm chứng toàn matrix. Nhóm này không đo tốc độ.

\section{Nhóm 2: ablation Taylor, prior và score}
\label{sec:g2}
\textbf{Setting chính:} cùng TRAIN pilot 224 cặp; giữ hai bộ tham số nhóm 1.
Ba cost: affine, exact relative và exact journal. Với mỗi kết quả $P$,
tính ba score mà không giải lại: Eq.19-style, spatial $\langle P,D\rangle$
và affine $\langle P,D+\mu R\rangle$. Bảng chính dùng residual stop
$r\leq10^{-3}$, tối đa 4000 vòng. Affine chạy Flash FP32; exact chạy dense GPU FP32.
\begin{center}\small
\begin{tabular}{@{}lrrr@{}}
\toprule
Model & Potentials MAP/ACC\% & Spatial MAP/ACC\% & Affine MAP/ACC\% \\
\midrule
Affine default & 63.770 / 57.143 & 54.924 / 35.714 & 64.290 / 57.143 \\
Affine tuned & 72.058 / 64.286 & 73.214 / 64.286 & 72.692 / 64.286 \\
Exact relative default & 64.335 / 57.143 & 54.860 / 35.714 & 64.335 / 57.143 \\
Exact relative tuned & 72.058 / 64.286 & 73.214 / 64.286 & 72.692 / 64.286 \\
Exact journal preset & 60.156 / 42.857 & 77.289 / 71.429 & 77.289 / 71.429 \\
Exact journal, sigma Flash & 50.950 / 42.857 & 50.950 / 42.857 & 50.950 / 42.857 \\
\bottomrule
\end{tabular}
\end{center}
Mỗi ô là MAP\% / ACC@1\%. Cả 224 cặp của sáu model đạt ngưỡng,
không có capped pair. Kết quả đã được kiểm tra lại; ranking CPU/GPU trùng
trong toàn bộ 54 phép so lưu của ablation.

Taylor/exact relative ở cấu hình selected giữ cùng MAP/ACC cho cả ba score,
dù $P$ khác tối đa khoảng 1.040\%; không suy ra Taylor vô hại trên mọi dữ liệu.
Giữ cùng $P$ của journal preset, đổi Eq.19-style sang spatial tăng ACC
từ 42.857\% lên 71.429\%, cho thấy score ảnh hưởng rõ đến kết quả.

\textbf{Giới hạn phép so:} ``sigma Flash'' là sigma chọn cho Flash relative
được đưa sang journal, không phải tuning riêng cho journal. Khác đơn vị prior
làm cùng giá trị sigma có độ mạnh khác nhau; kết quả thấp của hàng này
không chứng minh journal kém hơn sau tuning. Eq.19-style ở exact chỉ là
score đối chứng; affine score không phải total cost journal.
Nhóm 2 là ablation TRAIN pilot, không phải benchmark tốc độ hay so TEST cuối cùng.

\section{Nhóm 3: tuning riêng từng phương pháp trên TRAIN}
\label{sec:g3}
\textbf{Setting chính:} ba stratified TRAIN holdout, mỗi split 28 gallery/28 query
không giao nhau. Bảy phương pháp có tuning thử sáu candidate mỗi phương pháp;
DTW/LDTW/NDTW/OT giữ một cấu hình. Hoàn tất 138 job, 108192 cặp.
Chọn mean ACC@1, MAP phá hòa; candidate có cặp chưa đạt residual $10^{-3}$
tại cap 4000 bị loại. Mỗi phương pháp giữ score riêng, không đổi theo TEST.
Flash/dense entropic dùng GPU FP32; DTW/Soft-DTW/OT dùng CPU FP64.
% BEGIN GENERATED GROUP3 RESULTS
\begingroup\small
\begin{longtable}{@{}lrrrr@{}}
\caption{Mean phần trăm (sample SD) trên ba split TRAIN; chưa phải TEST.}\label{tab:g3results}\\
\toprule
Metric & Preset ACC@1 & Preset MAP & Selected ACC@1 & Selected MAP \\
\midrule
\endfirsthead
\toprule
Metric & Preset ACC@1 & Preset MAP & Selected ACC@1 & Selected MAP \\
\midrule
\endhead
\bottomrule
\endfoot
FlashOPW main & 58.333 (5.455) & 58.464 (5.594) & 79.762 (14.434) & 74.405 (7.329) \\
DTW & 67.857 (7.143) & 65.688 (7.305) & 67.857 (7.143) & 65.688 (7.305) \\
LDTW & 67.857 (7.143) & 65.688 (7.305) & 67.857 (7.143) & 65.688 (7.305) \\
NDTW & 64.286 (6.186) & 61.321 (6.824) & 64.286 (6.186) & 61.321 (6.824) \\
Soft-DTW & 72.619 (7.435) & 69.106 (7.646) & 78.571 (7.143) & 72.956 (5.695) \\
OT & 41.667 (4.124) & 47.484 (3.464) & 41.667 (4.124) & 47.484 (3.464) \\
Sinkhorn & 39.286 (3.571) & 47.527 (3.019) & 41.667 (2.062) & 47.268 (4.137) \\
TLp & 80.952 (12.542) & 75.130 (8.109) & 80.952 (12.542) & 75.130 (8.109) \\
OPW-KL & 77.381 (12.542)$^{*}$ & 72.899 (8.192)$^{*}$ & 76.190 (13.521) & 71.376 (7.876) \\
TCOT & 35.714 (6.186) & 44.922 (4.825) & 46.429 (9.449) & 49.580 (5.957) \\
OPW journal & 77.381 (12.542)$^{*}$ & 72.899 (8.192)$^{*}$ & 76.190 (13.521) & 71.397 (7.897) \\
\end{longtable}
\endgroup
$^{*}$ Preset không đủ điều kiện freeze vì có cặp chưa đạt tau tại cap.

\begingroup\small
\begin{longtable}{@{}lp{.68\linewidth}@{}}
\caption{Bộ tham số selected đã freeze; mọi metric cost scale bằng 1.}\label{tab:selected}\\
\toprule
Metric & Tham số \\
\midrule
\endfirsthead
\toprule
Metric & Tham số \\
\midrule
\endhead
\bottomrule
\endfoot
FlashOPW main & $\lambda_1=1$, $\lambda_2=0.1$, $\sigma=0.03194382825$ \\
DTW & Không có tham số tuning trong protocol này \\
LDTW & Không có tham số tuning trong protocol này \\
NDTW & Không có tham số tuning trong protocol này \\
Soft-DTW & $\gamma=1$ \\
OT & Không có tham số tuning trong protocol này \\
Sinkhorn & $\varepsilon=0.01$ \\
TLp & $w=50$, $\varepsilon=0.1$ \\
OPW-KL & $\lambda_1=0$, $\lambda_2=0.3$, $\sigma=5$ \\
TCOT & $\lambda=3$ \\
OPW journal & $\lambda_1=1$, $\lambda_2=0.3$, $\sigma=5$ \\
\end{longtable}
\endgroup
% END GENERATED GROUP3 RESULTS
Selected entropic đạt ngưỡng trên toàn bộ 2352 cặp mỗi phương pháp;
kết quả đã được kiểm tra lại từ matrices. FP64 verification trên 36 score
selected có max absolute error $6.609\cdot10^{-5}$, riêng Flash
$1.023\cdot10^{-6}$.

Flash tăng 21.429 điểm ACC và 15.941 điểm MAP so preset trên cùng protocol.
TLp cao hơn Flash 1.190 điểm ACC và 0.725 điểm MAP, tương ứng 68 so với
67 dự đoán đúng trong 84 lượt validation. Ba split chỉ có 71 query khác nhau;
độ biến động lớn, chưa kết luận ưu thế TEST hoặc statistical superiority.
Ngân sách bằng số attempt, không bằng runtime; grid nhỏ không bảo đảm global optimum.

\subsection{Đối chiếu FlashOPW main với OPW journal}
Cùng ba TRAIN split và kết quả nhóm 3; mỗi phương pháp dùng tham số selected
riêng và native score, không phải một lượt solve mới.
\begin{center}\small
\begin{tabular}{@{}llrrr@{}}
\toprule
Profile & Method & ACC@1\% (SD) & MAP\% (SD) & Cặp vượt tau \\
\midrule
Preset & Flash main & 58.333 (5.455) & 58.464 (5.594) & 0 \\
Preset & OPW journal & 77.381 (12.542) & 72.899 (8.192) & 1 \\
Selected & Flash main & 79.762 (14.434) & 74.405 (7.329) & 0 \\
Selected & OPW journal & 76.190 (13.521) & 71.397 (7.897) & 0 \\
\bottomrule
\end{tabular}
\end{center}
Sau tuning riêng, mean TRAIN của Flash cao hơn journal trong phép so này.
Đây là kết quả sau search trên TRAIN, chưa chứng minh Flash tốt hơn trên official TEST.

\section{Nhóm 4: official TEST --- chưa có kết quả đầy đủ}
\label{sec:g4}
\textbf{Setting chính:} toàn TRAIN gallery 200, TEST query 2050;
11 phương pháp, preset và selected riêng đã freeze từ TRAIN,
17 cấu hình duy nhất sau gộp cấu hình trùng. Mỗi cấu hình 410000 cặp.
Giữ residual $10^{-3}$, cap 4000; ACC@1 chính, MAP toàn gallery.
Phân tích dự kiến dùng paired query bootstrap và McNemar với Holm adjustment.

Trong dữ liệu đã nhận, full TEST chưa hoàn tất và chưa có bảng kết quả
cuối được kiểm tra đầy đủ. Không trình bày ACC/MAP/CI dự đoán hoặc kết quả
một phần như kết quả cuối cùng.

\section{Nhóm 5: speed/memory --- kết quả CPU sơ bộ}
\label{sec:g5}
\textbf{Setting chính:} synthetic sin trajectories với time warp và noise;
$N=512,1024,2048$, $M=1.25N$, $d=1,13$, hai seed.
Bốn phương pháp DTW/Soft-DTW/Sinkhorn/OPW, 48 case CPU FP64,
1 warmup và 2 repeats. Sinkhorn 100 vòng; OPW 20 vòng với $(1,0.1,1)$.
Deadline 45 s cho toàn process gồm startup, warmup và repeats.

\textbf{Kết quả cuối:} 37/48 case hoàn tất, 11 timeout, không code failure;
timing chi tiết ở Phụ lục~\ref{app:cpu-scaling}. Đây là kiểm tra chức năng/scaling,
có tải CPU chia sẻ chưa kiểm soát. Chưa có same-GPU benchmark để kết luận
speedup hoặc lợi thế memory FlashOPW; timeout không được dùng làm timing giả.

\section{Nhóm 6: k-NN end-to-end --- kiểm tra pipeline}
\label{sec:g6}
\textbf{Setting chính:} tiny synthetic dataset, 6 gallery/4 query, $d=2$;
CPU FP64, 5 vòng affine, k=1/3; kiểm tra default/selected và dense controls.
\textbf{Kết quả cuối:} integration không failure; max score error
$7.11\cdot10^{-15}$ với default, $8.88\cdot10^{-16}$ với selected,
NN agreement 100\%. Kết quả chỉ xác nhận pipeline trên tiny dataset.
Chưa có timing end-to-end đầy đủ cho transfer, distance matrix, ranking
và tổng thời gian trên workload thực với tham số frozen.

\section{Kết luận}
FlashSinkhorn tự triển khai có bằng chứng correctness và tám benchmark GPU,
nhưng chưa tái hiện toàn bộ speedup/thứ hạng của paper A100.
FlashOPW có parity tốt với reference cùng công thức main trên pilot;
ablation cho thấy prior và score cần được tách khỏi ảnh hưởng Taylor.
Tuning riêng trên TRAIN cải thiện Flash so preset, nhưng chưa đủ kết luận
ưu thế official TEST. Nhóm 4 cần kết quả full TEST; nhóm 5--6 cần benchmark
GPU được kiểm soát tải và đo toàn pipeline.

\appendix
\clearpage
\section{Toàn bộ số đo FlashSinkhorn}
\label{app:fs-tables}
Số liệu giữ precision console; backward nghĩa là forward cộng backward.
Memory là peak PyTorch allocated gồm live tensors, không tổng VRAM.
\SKIP\ là preflight skip; \NR\ là không được báo cáo, không phải số đo bằng 0.
% BEGIN GENERATED FLASH TABLES
\begingroup\small
\begin{longtable}{@{}rrrrrrr@{}}
\caption{Experiment 1: forward theo point count, d=64 (ms).}\label{tab:forward-n}\\
\toprule
$n=m$ & $d$ & Flash-sym & Flash-alt & KeOps & Tensorized & JAX \\
\midrule
\endfirsthead
\toprule
$n=m$ & $d$ & Flash-sym & Flash-alt & KeOps & Tensorized & JAX \\
\midrule
\endhead
\bottomrule
\endfoot
5000 & 64 & 2.259 & 2.182 & 42.87 & 25.104 & 8.795 \\
10000 & 64 & 6.613 & 7.64 & 76.98 & 105.038 & 25.182 \\
20000 & 64 & 24.337 & 26.502 & 273.124 & 422.083 & 82.611 \\
30000 & 64 & 49.996 & 54.45 & 404.046 & \SKIP & 177.591 \\
40000 & 64 & 88.538 & 89.333 & 796.254 & \SKIP & 307.218 \\
50000 & 64 & 135.652 & 138.441 & 1311.162 & \SKIP & 488.954 \\
\end{longtable}
\endgroup

\begingroup\small
\begin{longtable}{@{}rrrrrrr@{}}
\caption{Experiment 2: forward theo dimension, n=m=20000 (ms).}\label{tab:forward-d}\\
\toprule
$n=m$ & $d$ & Flash-sym & Flash-alt & KeOps & Tensorized & JAX \\
\midrule
\endfirsthead
\toprule
$n=m$ & $d$ & Flash-sym & Flash-alt & KeOps & Tensorized & JAX \\
\midrule
\endhead
\bottomrule
\endfoot
20000 & 4 & 13.477 & 11.842 & 30.639 & 422.091 & 70.41 \\
20000 & 8 & 13.455 & 11.806 & 44.369 & 422.065 & 71.62 \\
20000 & 16 & 13.545 & 11.81 & 76.857 & 422.093 & 73.263 \\
20000 & 32 & 13.508 & 11.85 & 143.585 & 422.075 & 76.313 \\
20000 & 64 & 24.503 & 26.387 & 271.441 & 422.069 & 82.223 \\
20000 & 128 & 45.02 & 47.333 & 437.425 & 422.071 & 103.432 \\
20000 & 256 & 120.886 & 126.721 & 1216.58 & 422.093 & 143.678 \\
20000 & 512 & 238.024 & 240.36 & 2354.527 & 422.098 & 231.932 \\
20000 & 1024 & 468.222 & 473.183 & 4666.855 & 422.079 & 476.648 \\
\end{longtable}
\endgroup

\begingroup\small
\begin{longtable}{@{}rrrrrrr@{}}
\caption{Experiment 3: forward cộng backward theo point count, d=64 (ms).}\label{tab:backward-n}\\
\toprule
$n=m$ & $d$ & Flash-sym & Flash-alt & KeOps & Tensorized & JAX \\
\midrule
\endfirsthead
\toprule
$n=m$ & $d$ & Flash-sym & Flash-alt & KeOps & Tensorized & JAX \\
\midrule
\endhead
\bottomrule
\endfoot
5000 & 64 & 3.835 & 3.254 & 50.804 & 31.924 & 9.188 \\
10000 & 64 & 8.21 & 9.212 & 90.313 & 132.29 & 29.122 \\
20000 & 64 & 28.502 & 30.628 & 315.469 & 530.739 & 93.925 \\
30000 & 64 & 58.154 & 62.554 & 469.833 & \SKIP & 204.235 \\
40000 & 64 & 101.63 & 102.717 & 916.931 & \SKIP & 357.857 \\
50000 & 64 & 155.901 & 160.041 & 1508.521 & \SKIP & 644.525 \\
\end{longtable}
\endgroup

\begingroup\small
\begin{longtable}{@{}rrrrrrr@{}}
\caption{Experiment 4: forward cộng backward theo dimension, n=m=20000 (ms).}\label{tab:backward-d}\\
\toprule
$n=m$ & $d$ & Flash-sym & Flash-alt & KeOps & Tensorized & JAX \\
\midrule
\endfirsthead
\toprule
$n=m$ & $d$ & Flash-sym & Flash-alt & KeOps & Tensorized & JAX \\
\midrule
\endhead
\bottomrule
\endfoot
20000 & 4 & 15.632 & 13.957 & 36.108 & 530.226 & 81.261 \\
20000 & 8 & 15.596 & 13.926 & 51.648 & 530.158 & 80.077 \\
20000 & 16 & 15.573 & 13.92 & 88.26 & 530.168 & 81.135 \\
20000 & 32 & 15.602 & 14.097 & 165.044 & 530.321 & 85.342 \\
20000 & 64 & 28.051 & 30.061 & 311.295 & 530.337 & 92.153 \\
20000 & 128 & 51.968 & 54.245 & 579.294 & 530.987 & 114.291 \\
20000 & 256 & 135.281 & 141.095 & 2683.626 & 534.058 & 163.367 \\
20000 & 512 & 273.385 & 275.735 & 17091.415 & 541.533 & 262.419 \\
20000 & 1024 & 538.959 & 544.009 & 71404.812 & 556.28 & 527.267 \\
\end{longtable}
\endgroup

\begingroup\small
\begin{longtable}{@{}rrrr@{}}
\caption{Experiment 5: peak PyTorch allocated forward, d=1024 (MB theo log).}\label{tab:memory-forward}\\
\toprule
$n=m$ & $d$ & Flash-alt & Tensorized \\
\midrule
\endfirsthead
\toprule
$n=m$ & $d$ & Flash-alt & Tensorized \\
\midrule
\endhead
\bottomrule
\endfoot
5000 & 1024 & 95.314 & 428.307 \\
10000 & 1024 & 174.08 & 1370.148 \\
20000 & 1024 & 325.714 & 5187.666 \\
30000 & 1024 & 482.607 & \SKIP \\
40000 & 1024 & 634.388 & \SKIP \\
50000 & 1024 & 787.202 & \SKIP \\
\end{longtable}
\endgroup

\begingroup\small
\begin{longtable}{@{}rrrr@{}}
\caption{Experiment 6: peak PyTorch allocated forward cộng backward, d=1024 (MB theo log).}\label{tab:memory-backward}\\
\toprule
$n=m$ & $d$ & Flash-alt & Tensorized \\
\midrule
\endfirsthead
\toprule
$n=m$ & $d$ & Flash-alt & Tensorized \\
\midrule
\endhead
\bottomrule
\endfoot
5000 & 1024 & 101.008 & 562.464 \\
10000 & 1024 & 184.974 & 2104.03 \\
20000 & 1024 & 345.042 & 8182.158 \\
30000 & 1024 & 512.45 & \SKIP \\
40000 & 1024 & 673.041 & \SKIP \\
50000 & 1024 & 839.925 & \SKIP \\
\end{longtable}
\endgroup

\begingroup\small
\begin{longtable}{@{}rrrrr@{}}
\caption{Experiment 7: HVP theo point count, d=64 (ms).}\label{tab:hvp-n}\\
\toprule
$n=m$ & $d$ & Flash-sym & KeOps & JAX \\
\midrule
\endfirsthead
\toprule
$n=m$ & $d$ & Flash-sym & KeOps & JAX \\
\midrule
\endhead
\bottomrule
\endfoot
5000 & 64 & 65.31 & 193.927 & 290.295 \\
6000 & 64 & 100.726 & 216.518 & 409.974 \\
7000 & 64 & 127.179 & 237.857 & 566.703 \\
8000 & 64 & 153.503 & 264.787 & 698.838 \\
9000 & 64 & 199.077 & 292.444 & 1056.555 \\
10000 & 64 & 236.438 & 315.049 & 1289.244 \\
20000 & 64 & 912.952 & 931.432 & \NR \\
30000 & 64 & 2032.442 & 1454.346 & \NR \\
40000 & 64 & 3532.364 & 2745.119 & \NR \\
50000 & 64 & 5337.679 & 4412.484 & \NR \\
\end{longtable}
\endgroup

\begingroup\small
\begin{longtable}{@{}rrrrr@{}}
\caption{Experiment 8: HVP theo dimension, n=m=10000 (ms).}\label{tab:hvp-d}\\
\toprule
$n=m$ & $d$ & Flash-sym & KeOps & JAX \\
\midrule
\endfirsthead
\toprule
$n=m$ & $d$ & Flash-sym & KeOps & JAX \\
\midrule
\endhead
\bottomrule
\endfoot
10000 & 4 & 77.987 & 124.908 & 151.746 \\
10000 & 8 & 78.009 & 133.39 & 169.775 \\
10000 & 16 & 78.061 & 159.214 & 230.08 \\
10000 & 32 & 78.56 & 207.22 & 450.829 \\
10000 & 64 & 236.456 & 314.645 & 1289.259 \\
10000 & 128 & 619.745 & 761.873 & 4528.709 \\
10000 & 256 & 972.212 & \NR & \NR \\
10000 & 512 & 2107.318 & \NR & \NR \\
\end{longtable}
\endgroup
% END GENERATED FLASH TABLES

\section{Số đo CPU scaling FlashOPW}
\label{app:cpu-scaling}
Median ms của hai repeats. TO là timeout 45 s cho toàn process,
không phải latency một phép toán. Tải CPU chưa được kiểm soát;
không dùng bảng này suy ra speedup Flash/GPU.
% BEGIN GENERATED CPU SCALING
\begingroup\small
\begin{longtable}{@{}rrrrrrrr@{}}
\caption{CPU scaling: median ms; TO = timeout.}\label{tab:cpu-scaling}\\
\toprule
$N$ & $M$ & $d$ & Seed & DTW & Soft-DTW & Sinkhorn & OPW \\
\midrule\endfirsthead
\toprule
$N$ & $M$ & $d$ & Seed & DTW & Soft-DTW & Sinkhorn & OPW \\
\midrule\endhead
\bottomrule\endfoot
512 & 640 & 1 & 42 & 3.714 & 20.042 & 3575.492 & 1974.366 \\
512 & 640 & 1 & 43 & 11.636 & 53.967 & 5493.932 & 1677.869 \\
512 & 640 & 13 & 42 & 10.550 & 32.628 & 6686.296 & 2424.992 \\
512 & 640 & 13 & 43 & 4.929 & 20.112 & 4698.369 & 1245.474 \\
1024 & 1280 & 1 & 42 & 27.023 & 89.335 & TO & 5794.937 \\
1024 & 1280 & 1 & 43 & 12.667 & 97.638 & 13143.665 & 4024.380 \\
1024 & 1280 & 13 & 42 & 35.884 & 151.147 & TO & 5578.916 \\
1024 & 1280 & 13 & 43 & 32.729 & 147.522 & TO & 6315.981 \\
2048 & 2560 & 1 & 42 & 79.946 & 535.837 & TO & TO \\
2048 & 2560 & 1 & 43 & 56.821 & 345.675 & TO & TO \\
2048 & 2560 & 13 & 42 & 111.559 & 475.469 & TO & TO \\
2048 & 2560 & 13 & 43 & 120.620 & 871.089 & TO & TO \\
\end{longtable}
\endgroup
% END GENERATED CPU SCALING

\begin{thebibliography}{9}
\bibitem{flashpaper}
Felix X.-F. Ye, Xingjie Li, An Yu, Ming-Ching Chang, Linsong Chu,
Davis Wertheimer.
\emph{FlashSinkhorn: IO-Aware Entropic Optimal Transport on GPU}.
arXiv:2602.03067v3, 2026.
\url{https://arxiv.org/abs/2602.03067v3}.

\bibitem{mainpdf}
Tài liệu phương pháp FlashOPW do người dùng cung cấp:
\code{main (2).pdf}, bản ngày 02/10/2026.

\bibitem{journalpdf}
Tài liệu OPW journal do người dùng cung cấp: \code{OWD_journal.pdf}.
\end{thebibliography}
\end{document}
